\documentclass[11pt,a4paper]{article}

\usepackage[margin=2.5cm]{geometry}
\usepackage{amsmath, amssymb, amsthm}
\usepackage{graphicx}
\usepackage{booktabs}
\usepackage{caption}
\usepackage{xcolor}
\usepackage[numbers,sort&compress]{natbib}
\usepackage[colorlinks=true, allcolors=blue]{hyperref}
\usepackage{cleveref}

\newtheorem{theorem}{Theorem}
\newtheorem{lemma}[theorem]{Lemma}

\title{A Concise Title That Describes Your Contribution}
\author{First Author\thanks{University of Somewhere, \texttt{first@example.org}}
   \and Second Author}
\date{}

\begin{document}
\maketitle

\begin{abstract}
One paragraph summarising the problem, the approach, the main result and why
it matters. Keep it under 200 words.
\end{abstract}

\section{Introduction}\label{sec:intro}
Motivate the problem and state your contributions. Prior work includes the
classic texts of \citet{knuth1984} and \citet{lamport1994}.

\section{Method}\label{sec:method}
\begin{theorem}\label{thm:main}
For every $n \in \mathbb{N}$, $\sum_{k=1}^{n} k = \frac{n(n+1)}{2}$.
\end{theorem}
\begin{proof}
By induction on $n$.
\end{proof}

\Cref{thm:main} is used throughout \cref{sec:results}.

\section{Results}\label{sec:results}
\begin{figure}[t]
  \centering
  \includegraphics[width=0.6\linewidth]{figures/results.png}
  \caption{Main result.}\label{fig:results}
\end{figure}

\begin{table}[t]
  \centering
  \caption{Comparison with baselines.}\label{tab:comparison}
  \begin{tabular}{lrr}
    \toprule
    Method & Accuracy (\%) & Time (s) \\
    \midrule
    Baseline & 81.2 & 12.4 \\
    Ours     & \textbf{89.7} & \textbf{7.9} \\
    \bottomrule
  \end{tabular}
\end{table}

As shown in \cref{fig:results} and \cref{tab:comparison}, our method wins.

\section{Conclusion}
Summarise and outline future work.

\bibliographystyle{plainnat}
\bibliography{references}
\end{document}
